%=============================================================================!
% 17.7  A PROTOCOL FOR A PRODUCTION RUN                                       !
%=============================================================================!
\section{A protocol for a production run}
\label{sec:ip-protocol}

A pilot runs through a checklist before every flight because missing one
item is costly, and the list keeps the order and the checks the same each
time. A production run of molecular
dynamics has such a list, and every item on it has been measured in an
earlier chapter. This section gathers them in the order a run meets them,
with the measurement behind each and what a healthy run reads.
\Cref{tab:ip-healthy} collects the readings.

\subsection*{Before the run}

The starting structure is checked before any step is taken: the closest
approach of two atoms, the largest force and the total momentum, which
should be zero (\cref{sec:cs-trajectory}). A structure built by hand, or
with a new potential, is first relaxed (\cref{sec:cv-protocol}), so that the
first steps do not turn stored strain into heat. Lattice constants come
from measurement where they exist, and the model's own relaxed values are
reported beside them, as \cref{sec:ip-learned,sec:ip-lic6} do for
graphite and LiC$_6$.
A learned potential is checked on the system before it is trusted:
against first-principles forces on configurations like those of the run
(\cref{sec:ip-aimd}), and for any physics its training data lacked, such
as the dispersion that holds graphite's sheets at their measured
spacing (\cref{sec:ip-learned}).

\subsection*{The time step}

The step is measured (\cref{sec:in-stability}). Short runs at fixed
energy with two or three steps record the fluctuation of the total
energy against that of the kinetic energy, which should be small and fall
in proportion to $\dt^2$ as the step is shortened, and the drift, which
should pass the test of \cref{sec:cv-drift} at the length of the
production run. The fastest vibration sets the limit: bonds to hydrogen
hold flexible water to \SI{0.5}{\femto\second}, its energy spreading by
6.7\% of the kinetic energy's at \SI{1}{\femto\second} (\cref{sec:cs-water}),
and the stiff
carbon-carbon bonds of graphite to the \SI{2}{\femto\second} that
\cref{sec:ip-learned} measures. A learned potential evaluated in single
precision adds an error of its own to the forces, and the same test, run
in both precisions, shows whether it matters.

\subsection*{Reaching the state}

A run is brought to its temperature and pressure first, and only then
sampled. Velocity rescaling and Berendsen coupling bring a system near a
temperature quickly, but their fluctuations are wrong and a long run under
them feeds slow motions (\cref{sec:th-berendsen}); Berendsen's barostat
likewise serves only to approach a pressure
(\cref{sec:pr-rescaling,sec:pr-trajectory}). The part of the run before
the averages settle is found by the search for a start of
\cref{sec:cv-equilibration} and discarded.

\subsection*{Sampling}

For averages over arrangements, any thermostat that passes the test of
\cref{sec:cv-ensemble} will do: CSVR passed it with $\tau_{\mathrm T} =
\SI{1}{\pico\second}$ for the liquid, Nosé-Hoover chains give the
canonical spread of the temperature (\cref{sec:th-compare}), and Langevin
dynamics samples the canonical ensemble even for a single atom
(\cref{sec:th-langevin}). For anything that depends on the motion, a
diffusion coefficient, a rate, a spectrum or a viscosity, the run is made
at fixed energy or under a global thermostat with a long time constant:
CSVR at \SI{1}{\pico\second} left the liquid's $D$ unchanged within its
error of 1\% (\cref{sec:cv-observables}), while a Langevin friction of
\SI{1}{\per\pico\second} slowed it by more than a quarter
(\cref{sec:th-langevin}). For constant pressure, stochastic cell rescaling
or a piston with chains samples the volume's fluctuations correctly, at
$\tau_{\mathrm P}$ near \SI{1}{\pico\second} for the liquid
(\cref{sec:pr-rescaling,sec:pr-piston}); a solid whose directions differ,
such as a layered one, needs its lengths coupled separately
(\cref{sec:pr-shape,sec:ip-lic6}).

\subsection*{How long}

The length of a run follows from the precision asked of each quantity:
$t_{\mathrm f} \ge 2\tau_{\mathrm{int}}\sigma_A^2/\sigma^{\ast2}$
(\cref{eq:cv-length}), with $\tau_{\mathrm{int}}$ and $\sigma_A$ measured
in a pilot run. For liquid argon under CSVR this ranged from
\SI{0.33}{\nano\second} for the pressure to \SI{10}{\nano\second} for the
temperature at the precisions of \cref{tab:cv-observables}. Each reported
average passes \texttt{convergence\_report}: a start in the first half of
the run, at least 100 effective samples, block errors that agree with the
correlation estimate, and no drift. A rate is reported only when enough
events are counted to give it an error bar, otherwise as unresolved, with
the run that would resolve it (\cref{sec:ip-hops}); a free-energy
difference is known only to about $2\kB T/\sqrt n$ from $n$ crossings of
the barrier (\cref{sec:ob-fes}).

\subsection*{After the run}

Transport coefficients from periodic boxes are corrected for the size of
the box, or measured at two sizes (\cref{sec:cv-size}). Displacements
are computed from unwrapped positions, removing centre-of-mass motion
where appropriate. Each observable is checked against an identity where one exists, such
as the energy from $g(r)$ or $D$ from both the MSD and the velocities
(\cref{sec:ob-trajectory}). Finally the record is passed through
\texttt{health\_report} (\cref{sec:ip-trajectory}), and the run is
reported with what makes it repeatable: the code and its version, the
model and its weights, every option of the integrator and thermostat with
its units, the seeds, and the error bar of every number.

\begin{table}[htb]
\centering
\small
\caption{What a healthy run reads, and where the book measured it.}
\label{tab:ip-healthy}
\begin{tabular}{@{}>{\raggedright\arraybackslash}p{0.29\linewidth}
   >{\raggedright\arraybackslash}p{0.47\linewidth}
   >{\raggedright\arraybackslash}p{0.16\linewidth}@{}}
\toprule
quantity & healthy reading & measured in\\
\midrule
total energy at fixed energy & spread small against $K$'s, falling as
$\dt^2$; no drift by the test of Ch.~15 & \cref{sec:in-stability,sec:cv-drift}\\
kinetic temperature & mean on target; spread $\sqrt{2/N_{\mathrm f}}$ of
it under a canonical thermostat, less at fixed energy &
\cref{sec:sm-kinetic,sec:th-tests}\\
velocities & excess kurtosis within $3\sqrt{24/3N}$ of zero &
\cref{sec:sm-maxwell}\\
centre of mass & still, or, under a thermostat that does not keep the
momentum, displacements measured from it &
\cref{sec:cs-trajectory,sec:th-trajectory}\\
volume under a barostat & spread $\sqrt{\kB T\ens{V}\kappa_T}$; every
coupled component of the pressure at $P_0$ & \cref{sec:pr-npt}\\
an average & passes \texttt{convergence\_report} & \cref{sec:cv-protocol}\\
positions for displacements & continuous, every move a small fraction of
the cell & \cref{sec:ob-trajectory}\\
a rate & at least 25 events for a fifth, or reported unresolved &
\cref{sec:ip-hops}\\
a learned potential & forces within the error measured against its
reference on this system & \cref{sec:ip-aimd}, still to be made\\
\bottomrule
\end{tabular}
\end{table}

\begin{keyidea}
A production run is relaxed, given a measured time step, brought to its
state with any thermostat and barostat, sampled with one that is
canonical, and, for dynamics, a global or no thermostat, run for the
length its target errors need, and checked against the readings of a
healthy run. Every choice in the list has a measurement behind it in this
book.
\end{keyidea}

\begin{selfcheck}
A diffusion coefficient is to be measured, and the only thermostat at
hand is Langevin. What would a careful run do?
\end{selfcheck}
